\chapter{Loop Effects in QED}

\begin{remarkbox}
Tree-level QED describes photon exchange, annihilation, and emission. Loop
corrections reveal the genuinely quantum structure of the theory: particles
interact with their own fields, the vacuum is polarizable, charges run with
scale, and infrared physics forces us to define inclusive observables.
\end{remarkbox}

\section{Why Loop Effects Matter}

The previous chapter introduced QED at tree level. Tree diagrams already explain
many qualitative features of electromagnetic scattering, but they are not the
whole theory. The perturbative expansion of QED contains loop diagrams, and
these diagrams have two important roles.

First, loops produce finite measurable corrections. Examples include the
anomalous magnetic moment of the electron and radiative corrections to
scattering cross sections. Second, loops produce ultraviolet divergences. These
are not ignored; they are absorbed into the definitions of physical masses,
fields, and charges through renormalization.

The main one-loop structures in QED are
\begin{itemize}
    \item the electron self-energy;
    \item the photon self-energy, also called vacuum polarization;
    \item the vertex correction;
    \item real and virtual soft-photon effects.
\end{itemize}
Together they show why QED is both highly predictive and conceptually subtle.

\section{Electron Self-Energy}

The electron self-energy is the correction to the electron propagator from the
emission and reabsorption of a virtual photon. The one-loop diagram has an
internal photon line and an internal electron line. Its contribution is written
schematically as
\[
    -i\Sigma(p)
    = (-ie)^2\int\frac{\dd^4k}{(2\pi)^4}\,
      \gamma^\mu
      \frac{i(\gamma^\rho(p_\rho-k_\rho)+m)}{(p-k)^2-m^2+i\epsilon}
      \gamma^\nu
      \frac{-i\eta_{\mu\nu}}{k^2+i\epsilon}
\]
in Feynman gauge.

The full electron propagator is then formally
\[
    S(p)
    = \frac{i}{\gamma^\mu p_\mu-m-\Sigma(p)}.
\]
The self-energy shifts the location and residue of the propagator pole. The pole
position is used to define the physical electron mass, while the residue is
related to field-strength renormalization.

\section{Mass and Wave-Function Renormalization}

The bare mass appearing in the Lagrangian is not directly the measured mass. The
measured mass is associated with the pole of the full propagator. One imposes a
renormalization condition such as
\[
    \gamma^\mu p_\mu-m-\Sigma(p)=0
    \qquad \text{at} \qquad p^2=m_{\mathrm{phys}}^2.
\]
This condition fixes the relation between the bare mass and the physical mass.

The residue of the pole defines the wave-function renormalization. Near the
physical pole, the full propagator has the form
\[
    S(p)\sim \frac{iZ_2}{\gamma^\mu p_\mu-m_{\mathrm{phys}}}
    +\text{regular terms}.
\]
The constant \(Z_2\) is the electron field-strength renormalization. It tells us
how the field in the Lagrangian is normalized relative to the one-particle pole
of the full propagator.

\section{Vacuum Polarization}

The photon self-energy arises from a fermion loop with two external photon legs.
It is called vacuum polarization because it describes how virtual charged pairs
modify the propagation of the electromagnetic field.

The one-loop tensor has the structure
\[
    i\Pi^{\mu\nu}(q)
    = -(-ie)^2\int\frac{\dd^4k}{(2\pi)^4}\,
      \operatorname{tr}\left[
      \gamma^\mu\frac{i(\gamma^\rho k_\rho+m)}{k^2-m^2+i\epsilon}
      \gamma^\nu\frac{i(\gamma^\sigma(k_\sigma+q_\sigma)+m)}{(k+q)^2-m^2+i\epsilon}
      \right].
\]
The initial minus sign is the closed-fermion-loop sign.

Gauge invariance constrains this tensor. It must be transverse:
\[
    q_\mu\Pi^{\mu\nu}(q)=0.
\]
Therefore it has the form
\[
    \Pi^{\mu\nu}(q)
    = (q^\mu q^\nu-q^2\eta^{\mu\nu})\Pi(q^2).
\]
This transversality is a direct consequence of current conservation and the Ward
identity.

\section{Charge Renormalization and the Running Coupling}

Vacuum polarization modifies the photon propagator. In effect, the electric
charge measured at one momentum scale differs from the charge measured at
another scale. The physical picture is that the vacuum contains virtual charged
pairs that screen or antiscreen charges depending on the theory.

In QED, the effective electromagnetic coupling increases slowly with energy. One
writes a scale-dependent coupling
\[
    \alpha(\mu)=\frac{e(\mu)^2}{4\pi}.
\]
At leading order with one charged fermion, the beta function has the form
\[
    \beta(e)=\mu\frac{\dd e}{\dd\mu}
    = \frac{e^3}{12\pi^2}+\cdots.
\]
The positive sign means that the coupling grows at shorter distances.

This running is usually tiny at ordinary energies, but conceptually it is one of
the central discoveries of renormalized QFT: coupling constants depend on the
scale at which they are defined.

\section{Vertex Correction}

The QED vertex also receives a one-loop correction. Instead of the tree-level
factor
\[
    -ie\gamma^\mu,
\]
the full vertex has the form
\[
    -ie\Gamma^\mu(p',p),
\]
where
\[
    \Gamma^\mu(p',p)=\gamma^\mu+\Lambda^\mu(p',p).
\]
Here \(\Lambda^\mu\) denotes loop corrections.

The most general on-shell vertex for a spin-\(1/2\) particle can be decomposed
into form factors:
\[
    \bar u(p')\Gamma^\mu u(p)
    = \bar u(p')\left[
      \gamma^\mu F_1(q^2)
      +\frac{i\sigma^{\mu\nu}q_\nu}{2m}F_2(q^2)
      \right]u(p),
\]
where \(q=p'-p\). The function \(F_1\) is related to electric charge, while
\(F_2(0)\) gives the anomalous magnetic moment.

\section{The Anomalous Magnetic Moment}

For a Dirac particle at tree level, the gyromagnetic factor is
\[
    g=2.
\]
Loop corrections modify this. The anomalous magnetic moment is defined by
\[
    a_e=\frac{g-2}{2}.
\]
At one loop in QED,
\[
    a_e=\frac{\alpha}{2\pi}+\cdots.
\]
This finite prediction is one of the most famous successes of perturbative QFT.
It comes from the vertex correction after renormalization.

The important conceptual point is that loops are not merely divergent nuisances.
They also contain finite physical effects that can be measured with extraordinary
precision.

\section{Ward Identity}

Gauge invariance relates the electron self-energy and the vertex correction. The
Ward--Takahashi identity is
\[
    q_\mu\Gamma^\mu(p+q,p)
    = S^{-1}(p+q)-S^{-1}(p).
\]
In the limit \(q\to0\), this implies a relation between the vertex
renormalization and the electron wave-function renormalization. In common
notation,
\[
    Z_1=Z_2.
\]
Here \(Z_1\) renormalizes the electron--photon vertex, while \(Z_2\) renormalizes
the electron field.

This relation is a powerful consequence of gauge invariance. It ensures that
charge renormalization is controlled by the photon field renormalization rather
than by independent vertex and electron-field effects.

\section{Ultraviolet and Infrared Divergences}

QED loop integrals can diverge in two different regimes. Ultraviolet divergences
come from large loop momentum:
\[
    |k|\to\infty.
\]
They reflect sensitivity to short-distance physics and are handled by
renormalization.

Infrared divergences come from small photon momentum:
\[
    |k|\to0.
\]
They occur because the photon is massless. Soft virtual photons can produce
divergent loop corrections, while real photons of arbitrarily low energy can be
emitted without being experimentally resolved.

The physical treatment of these two kinds of divergences is different.
Ultraviolet divergences are absorbed into parameters. Infrared divergences
cancel only after defining observables inclusively enough.

\section{Soft Photon Emission}

Consider a process involving charged particles. In addition to the process with
no extra photons, there is a process with an additional soft photon:
\[
    \text{charged process}\longrightarrow \text{same process}+\gamma_{\mathrm{soft}}.
\]
If the photon energy is below detector resolution, the two outcomes cannot be
experimentally distinguished.

The amplitude for soft photon emission factorizes in a universal way. For
external charged particles with momenta \(p_i\), the soft factor has the form
\[
    e\sum_i \eta_i\frac{p_i\cdot\epsilon(k)}{p_i\cdot k},
\]
where \(\eta_i\) depends on whether the charged particle is incoming or outgoing
and on the sign of its charge.

The denominator \(p_i\cdot k\) becomes singular as the photon momentum \(k\) goes
to zero. This is the source of soft divergences in real emission.

\section{Inclusive Observables}

Virtual soft-photon loops and real soft-photon emission produce divergences with
opposite signs. Physical cross sections must include all experimentally
indistinguishable final states. Therefore one computes inclusive quantities such
as
\[
    \sigma_{\mathrm{inclusive}}
    = \sigma_{0}+\sigma_{\mathrm{virtual}}+\sigma_{\mathrm{soft\ real}}+\cdots.
\]
The infrared divergences cancel in this inclusive sum.

This cancellation is an example of a general principle: the mathematical object
that is finite is not always the exclusive amplitude for a fixed number of
massless particles. It is the observable that matches what a detector can
actually distinguish.

\section{Physical Meaning of Radiative Corrections}

Radiative corrections change how charged particles interact. They shift masses,
renormalize fields, modify the effective charge, correct magnetic moments, and
alter cross sections. They also teach us that the vacuum is not empty in a
classical sense: it supports virtual fluctuations that affect observable
processes.

In QED these effects are small because \(\alpha\) is small at low energies. This
makes QED one of the most precise theories in physics. The same logic applies in
other quantum field theories, although the coupling may be larger and the
calculations more difficult.

\begin{summarybox}
Loop effects in QED include electron self-energy, vacuum polarization, and vertex
corrections. They produce ultraviolet divergences that require renormalization
and finite effects such as the anomalous magnetic moment. Gauge invariance gives
Ward identities, including the relation \(Z_1=Z_2\). Because the photon is
massless, QED also has infrared divergences, which cancel only in inclusive
observables that include unresolved soft radiation.
\end{summarybox}

\section{Chapter Checklist}

After reading this chapter, one should be able to:
\begin{itemize}
    \item identify the main one-loop diagrams in QED;
    \item explain electron self-energy and its relation to mass renormalization;
    \item explain vacuum polarization and transversality;
    \item describe charge renormalization and the running coupling;
    \item write the form-factor decomposition of the QED vertex;
    \item state the one-loop result \(a_e=\alpha/(2\pi)+\cdots\);
    \item state the Ward--Takahashi identity and the implication \(Z_1=Z_2\);
    \item distinguish ultraviolet and infrared divergences;
    \item explain why real soft-photon emission must be included;
    \item explain why inclusive observables are infrared finite.
\end{itemize}
